
\section{Exact unforced generation of nonaxisymmetric Reynolds torque}
\label{sec:xxxvii-exact-unforced-torque}

The necessary transport condition of Edition~XXXV left open
whether an \emph{exact} unforced, reflection-breaking velocity
can have nonzero Reynolds angular-momentum production.
We answer this subproblem constructively on the actual
whole-space geometry $\mathbb R^3$. Our example is globally
smooth and deliberately does not preserve the forced
singularity of \cite{openai2026}; it must not be
mistaken for a solution of the unforced blowup problem.

Fix $\nu>0$, $r^2=x^2+y^2$, $\rho^2=r^2+z^2$,
and put $\varepsilon=\nu/1000$. Define a smooth Gaussian
vector potential and divergence-free datum by
\begin{equation}\label{eq:xxxvii-datum}
F(x,y,z)=\varepsilon e^{-\rho^2}
\bigl(x+r^2y\bigr),\qquad
u_0=\operatorname{curl}(F e_z)
=(\partial_yF,-\partial_xF,0).
\end{equation}
This datum breaks the $z$-axis rotational symmetry
and the reflection $(x,y,z)\mapsto(x,-y,z)$,
although it is smooth on the entire symmetry axis.

\begin{theorem}[A completely specified global zero-force flow]
\label{thm:xxxvii-global}
For the datum \eqref{eq:xxxvii-datum} there is
a unique global smooth finite-energy solution $u(t)$
of the exact incompressible \emph{unforced}
three-dimensional Navier--Stokes equations on $\mathbb R^3$:
\[
\partial_tu+(u\cdot\nabla)u-\nu\Delta u+\nabla p=0,
\qquad \operatorname{div}u=0,\qquad u(0)=u_0.
\]
The solution is defined without an external force
by the convergent projected mild fixed point, with
the physical pressure reconstructed from $u\otimes u$.
Using the unitary Fourier transform, set
$C_F=(2\pi)^{-3/2}$ and
$\|v\|_{\mathcal X^j}=\int|\xi|^j|\widehat v(\xi)|d\xi$.
Then the following explicit bounds hold:
\begin{align}
\boxed{\quad
q:=\frac{C_F}{\nu}\|u_0\|_{\mathcal X^{-1}}
&\le\frac{1}{100\sqrt\pi}<\frac1{16},
\qquad 16q\le\frac4{25\sqrt\pi}<1,}
\label{eq:xxxvii-critical}\\
\boxed{\quad
\sup_{t\ge0}\|u(t)\|_\infty
&\le\frac{33\nu}{2000}
\exp\left(\frac2{25\sqrt\pi}\right).}
\label{eq:xxxvii-uniform}
\end{align}
The global flow is nonaxisymmetric and
reflection-breaking for at least a positive
initial time interval.
\end{theorem}
\begin{proof}
The Fourier transform of the scalar
$e^{-\rho^2}(x+r^2y)$ is, with
$\xi=(\xi_x,\xi_y,\xi_z)$,
\[
\widehat{e^{-\rho^2}(x+r^2y)}
=\frac{i}{8\,2^{3/2}}\,
e^{-|\xi|^2/4}
\bigl((\xi_x^2+\xi_y^2)\xi_y
       -4\xi_x-8\xi_y\bigr).
\]
The modulus is dominated by
$2^{-3/2}e^{-|\xi|^2/4}
(|\xi|^3+12|\xi|)/8$.
Exact radial Gaussian integration gives
\[
\int_{\mathbb R^3}
\left|\widehat{e^{-\rho^2}(x+r^2y)}\right|d\xi
\le20\sqrt2\,\pi,\qquad
\int|\xi|
\left|\widehat{e^{-\rho^2}(x+r^2y)}\right|d\xi
\le33\sqrt2\,\pi^{3/2}.
\]
Since
$|\widehat{\operatorname{curl}(F e_z)}|
=|\xi_\perp||\widehat F|\le|\xi||\widehat F|$,
the two bounds yield
$q\le (10/\sqrt\pi)(\varepsilon/\nu)
=1/(100\sqrt\pi)$ and
$C_F\|u_0\|_{\mathcal X^0}
\le(33/2)\varepsilon$.
The global Fourier mild contraction in the
Banach norm
\[
\|u\|_{\mathscr B}
=\sup_{t\ge0}\|u(t)\|_{\mathcal X^{-1}}
+\nu\int_0^\infty\|u(t)\|_{\mathcal X^1}dt
\]
uses a ball of radius
$4\|u_0\|_{\mathcal X^{-1}}$
and has Lipschitz factor at most $16q<1$,
as calculated in Edition~XXXIV.
The data are Schwartz, so classical
parabolic smoothing gives a global smooth solution,
with no added force. The same Fourier estimate gives
$\|u(t)\|_{\mathcal X^0}\le
\|u_0\|_{\mathcal X^0}e^{8q}$,
proving \eqref{eq:xxxvii-uniform}.
Nonaxisymmetry and the broken reflection
hold for an initial time interval by smooth
continuity of $u(t)$ to $u_0$.
\end{proof}

\begin{theorem}[Strict nonzero Reynolds torque generated by the exact PDE]
\label{thm:xxxvii-exact-production}
For $u(t)$ from Theorem~\ref{thm:xxxvii-global},
write in cylindrical coordinates $\Gamma=ru_\theta$,
$\overline{g}=(2\pi)^{-1}\int_0^{2\pi}g\,d\theta$,
$g'=g-\bar g$ and
\[
\mathcal T=
-\frac1r\partial_r
(r\,\overline{u_r'\Gamma'})
-\partial_z\overline{u_z'\Gamma'}.
\]
Its initial mean angular momentum vanishes identically,
$\bar\Gamma(0,r,z)\equiv0$, but the initial
nonaxisymmetric flux and torque satisfy the
\emph{global exact formulas}
\begin{align}
\overline{u_r'(0)\Gamma'(0)}
&=\varepsilon^2r^3e^{-2(r^2+z^2)},\\
\boxed{\quad
\mathcal T(0,r,z)
&=4\varepsilon^2r^2(r^2-1)e^{-2(r^2+z^2)}.}
\label{eq:xxxvii-exact-torque}
\end{align}
At the fixed location $(r_*,z_*)=(1/2,0)$,
\begin{equation}\label{eq:xxxvii-positive-production}
\boxed{\quad
\partial_t\bar\Gamma(0,r_*,z_*)
=\mathcal T(0,r_*,z_*)
=-\frac34\varepsilon^2e^{-1/2}\ne0.}
\end{equation}
Consequently there exists $\delta>0$ such that
for all $0<t\le\delta$,
\begin{equation}\label{eq:xxxvii-positive-time}
\boxed{\quad
\mathcal T(t,r_*,z_*)\le
-\frac38\varepsilon^2e^{-1/2},\qquad
\bar\Gamma(t,r_*,z_*)\le
-\frac38\varepsilon^2e^{-1/2}t.}
\end{equation}
The location $r_*$, datum $\varepsilon$ and
initial production amplitude are numerically
specified. The positive existence interval $\delta$
is not given an independently interval-certified
numerical value.
\end{theorem}
\begin{proof}
Set $f(r,z)=e^{-(r^2+z^2)}$ and use
$F=\varepsilon f(r\cos\theta+r^3\sin\theta)$.
Taking the curl gives
\[
u_r(0)=\varepsilon f(r^2\cos\theta-\sin\theta),
\]
\[
u_\theta(0)=\varepsilon f\bigl[
(2r^2-1)\cos\theta+
(2r^4-3r^2)\sin\theta\bigr],\quad u_z(0)=0.
\]
The angular means of the sine/cosine components
of $\Gamma$ vanish. Direct angular integration gives
$\overline{u_r'\Gamma'}=\varepsilon^2r^3 f^2$
and $\overline{u_z'\Gamma'}=0$.
Radial differentiation proves
\eqref{eq:xxxvii-exact-torque}.
The exact angularly averaged projected
Navier--Stokes momentum equation was derived
in Edition~XXXV:
\[
(\partial_t+\bar u_r\partial_r+\bar u_z\partial_z)
\bar\Gamma
=\nu(\partial_{rr}-r^{-1}\partial_r+\partial_{zz})
\bar\Gamma+\mathcal T.
\]
Since $\bar\Gamma(0)\equiv0$, every
transport and diffusion term involving the
initial mean vanishes, so its initial
time derivative equals the displayed torque.
Equivalently, averaging the \emph{complete}
initial quadratic advection gives
\[
-\overline{u_r\partial_r\Gamma+
   (u_\theta/r)\partial_\theta\Gamma}
=4\varepsilon^2r^2(r^2-1)e^{-2(r^2+z^2)},
\]
an independent calculation from the full PDE.
At $r=1/2,z=0$ the result is negative.
Smooth PDE evolution makes $\mathcal T$
continuous locally in $(t,r,z)$; hence it
retains at least half this negative magnitude
for a short interval.
Integrate the mean evolution at $(r_*,z_*)$.
Its transport/diffusion terms are initially
zero but need not vanish at later times.
Nevertheless the mean time derivative itself
is continuous with the same strictly negative
initial value, so possibly shrinking $\delta$
makes $\partial_t\bar\Gamma\le
-3\varepsilon^2 e^{-1/2}/8$ there.
Integration gives the second inequality.
\end{proof}

\begin{proposition}[Why this exact torque generator does not preserve the forced singular core]
\label{prop:xxxvii-regularity-negative}
The globally smooth exact unforced solution from
Theorem~\ref{thm:xxxvii-global} obeys
\[
|u_\theta(x_\tau,1-\tau)|
\le\frac{33\nu}{2000}e^{2/(25\sqrt\pi)}
\quad(0<\tau<1)
\]
on the source's shrinking path
$x_\tau=(\sqrt{2X_{\rm in}\tau},0,0)$.
Therefore it cannot retain any positive fixed fraction of
$u_{B,\theta}(x_\tau,1-\tau)
\sim e_0\tau^{-1/2-h}$.
In the exact source correction $w=u-u_B$,
the growing source swirl cancels at leading order,
even though $\mathcal T$ is genuinely nonzero
for an initial time interval.
\end{proposition}
\begin{proof}
Use the uniform bound \eqref{eq:xxxvii-uniform}
and the original Section~10 core asymptotic.
\end{proof}

\paragraph{Proof boundary.}
The new point is an \emph{actual unforced nonlinear
production mechanism}, rather than an arbitrary
kinematic angular-momentum flux.
It proves only local-in-time nonzero mean torque.
Nothing here yields the necessary
$\tau^{-h}$ late-time nonaxisymmetric escape rate,
a singularity-preserving correction,
source-specific numerical cutoffs,
or an end-to-end Lean PDE certificate.
The global existence ingredient is classical
small-data theory, and the exact vortex formula
is a constructed test case, not a regularity breakthrough.
